\documentclass{article}
\usepackage[margin=2cm]{geometry}
\usepackage{amsmath}
\usepackage{graphicx}
\usepackage{booktabs}
\usepackage[utf8]{inputenc}
\usepackage{ragged2e}
\setlength{\emergencystretch}{3em}
\begin{document}
\sloppy

\noindent\hfill\includegraphics[height=1.8em]{assets/figure_header_logo.png}\par\noindent\hrule\vspace{4pt}

be of the form$V . { a - 1 } ^ { ~ a }$ as we see below, which means that
multiplication of HS generators we have to consider is


Inserting the normalisation $N . { * }$ we obtain


Then


\begin{equation}
\begin{array} { r } { \qquad }\end{array}
\end{equation}


\begin{equation}
\bar { V _ { - 1 } ^ { 2 } \star V _ { - 1 } ^ { 2 } \star.... \star V _ { - 1 } ^ { 2 } }. \qquad \qquad ( 1 )
\end{equation}


\begin{equation}
\begin{array} { c } { \phi _ { \bar { z } \dots \bar { z } } | _ { \partial \Lambda ^ { ( s ) } } = \frac { 1 } { 2 ^ { s - 1 } } e ^ { 2 ( s - 1 ) \rho } \bar { \partial } \Lambda ^ { ( s ) } } \\ { \times ~ \frac { 3 \cdot 4 \sqrt { \pi } 4 ^ { 4 - 2 s } \Gamma ( s ) \Gamma ( s + \lambda ) \Gamma ( s - \lambda ) } { ( \lambda ^ { 2 } - 1 ) \Gamma ( s + \frac { 1 } { 2 } ) \Gamma ( 1 - \lambda ) \Gamma ( 1 + \lambda ) }. \quad ( 1 1 ) } \end{array}
\end{equation}


where the$g . { 2 } ^ { 2 2 } ( - 1, - 1 ) ~ = ~ g . { 3 } ^ { 2 2 } ( - 1, - 1 ) ~ = ~ 0$ Multiplying
with following$V . { - 1 } ^ { 2 }$ etc. on e can conclude


The expression $\phi . { z, \dots, z }$ we want to compare with expres
sion ( for highest derivative on$C . { 0 }$ and $\partial \Lambda ^ { ( \pi ) }$ In the
computation of the vertex this would be a term


\begin{equation}
\underbrace { V _ { - 1 } ^ { 2 } * V _ { - 1 } ^ { 2 } *.... * V _ { - 1 } ^ { 2 } } _ { s - 1 } = \frac { 1 } { 2 ^ { s - 1 } } g _ { 1 } ^ { 2 ( s - 1 ) } ( - 1.. - ( s - 2 ) ) V _ { - ( s - 1 ) } ^ { s }
\end{equation}


\begin{equation}
\phi ^ { z... z } \phi \nabla _ { z }... \nabla _ { z } \phi \qquad ( 1 2 )
\end{equation}


while


\begin{equation}
\begin{array} { l } { g _ { 2 } ^ { n ~ s + m, s } } \\ { g _ { 2 } ^ { 2 ( s - 1 ) } ( - 1, - ( s - 2 ) ) = g _ { 3 } ^ { 2 ( s - 1 ) } ( - 1, - ( s - 2 ) ) = 0. ~ \mathrm { ~ ( 4 ) ~ } } \\ { \mathrm { T h a t ~ m e a r e ~ w u n ~ h a v e ~ f o m a n d ~ t h e ~ c r o n t r i b u t i o n ~ t h e ~ r o ~ t h e ~ } \quad}\end{array}
\end{equation}


for$\phi ^ { z \cdots z }$higher spin field with indices and$. { s }$  scalar field$\phi$
Raising indices contributes with a factor$2 ^ { s } e ^ { - 2 s \rho }$ so that
the field$\phi ^ { z \cdots z }$becomes


That means we have found the contribution to the z....
component multiplied with lowest derivative on$\mathbf { \Lambda } ^ { ( \pmb { s } ) }$due
to definition of trace for generators $\boldsymbol { v } . { \mathbf { n } } ^ { \star }$


\begin{equation}
\begin{array} { c } { \cdot ~ }\end{array}
\end{equation}


\begin{equation}
\begin{array} { l } { \phi ^ { z \dots z } = \frac { 1 } { 2 } e ^ { - 2 \rho } \bar { \partial } \Lambda ^ { ( s ) } \mathbf { 3 \cdot 4 ^ { 4 - 2 s } } } \\ { ~ }\end{array}
\end{equation}


and$\begin{array} { r } { \{ \mathbf { a } \} . { n } = \frac { \Gamma ( \mathbf { a } + n ) } { \Gamma ( \mathbf { a } ) } } \end{array}$ascending Pochhammer symbol. The
overall constant is set to.


\begin{equation}
N _ { \star } \equiv \frac { 3 \cdot 4 ^ { s - 3 } \sqrt { \pi } q ^ { 2 s - 4 } \Gamma ( s ) } { ( \lambda ^ { 2 } - 1 ) \Gamma ( s + \frac { 1 } { 2 } ) } ( 1 - \lambda ) _ { s - 1 } ( 1 + \lambda ) _ { s - 1 } \quad ( 6 )
\end{equation}


KG|highest number of derivatives(C)loA
$( - 1 ) ^ { * } 4 e ^ { - s p } \bar { \partial } \Lambda ^ { ( s ) } \partial ^ { * } C . { 0 } ^ { 1 }$we get (schematically written)


Let us go back to $\phi . { \xi \dots z }$ component.The star product
$k : 2, \cdots, k : 2$ will contribute with $\frac { 1 } { 2 ^ { n - 1 } } e ^ { ( n - 1 ) \rho } V . { - ( n - 1 ) } ^ { s }$if we
consider as explained above the lowest derivative on$\mathbf { \Lambda } ^ { ( \cdot ) }$
We can denote this as


\begin{equation}
t r ( V _ { 1 } ^ { 2 } V _ { - 1 } ^ { 2 } ) = - 1. \qquad \qquad \{ 7 \}
\end{equation}


\begin{equation}
\begin{array} { r l } & { \frac { \phi ^ { s... z } \big | _ { \tilde { G } \Lambda ^ { ( s ) } } } { \prod _ { K G } | _ { \mathrm { h i g h e x t ~ m u n b u r ~ o f e r i v e s } ( \delta C ) h } | _ { \tilde { G } \Lambda ^ { ( s ) } } } } \\ & { = ( - 1 ) ^ { a } \frac { 4 ^ { 4 - 2 s } \Gamma ( s ) \Gamma ( s + \lambda ) \Gamma ( s - \lambda ) } { ( \lambda ^ { 2 } - 1 ) \Gamma ( s + \frac {}}}\end{array}
\end{equation}


\begin{equation}
\begin{array} { l } { \mathrm { r e r o n s ~ c a n t e r e c o r ~ v i s e r ~ s u s } } \\ { \mathrm { ~ } e _ { s } *... * e _ { s } ( V _ { - 1 } ^ { 2 } *... * V _ { - 1 } ^ { 2 } ) = \frac { 1 } { 2 ^ { \pi - 1 } } e ^ { ( s - 1 ) \rho } V _ { - ( s - 1 ) } ^ { s }. ~ ~ ( 8 ) } \\ { \mathrm { ~ } ~ }\end{array}
\end{equation}


which taking into account the normalisation gives the
coupling for the 00s three point function.


\begin{equation}
\begin{array} { r } { r \times \cdots - \cdots - \cdots + \cdots } \\ { \partial _ { \tilde { z }, \dots, \tilde { z } } | _ { \partial \Lambda ^ { ( s ) } } = t r \left [ \frac { 1 } { 2 ^ { s - 1 } } e ^ { ( s - 1 ) } V _ { = ( x - 1 ) } ^ { s } \star V _ { s - 1 } ^ { s } e ^ { ( s - 1 ) \rho } \bar { \partial } \Lambda ^ { ( s ) } ( z, \tilde { z } ) \right ] } \\ { ~ }\end{array}
\end{equation}


The $\bar { E } . { \pmb { \xi } . { s } } ~ = ~ \bar { A } . { \pmb { \bar { z } } . { s } } ~ - ~ \bar { A } . { \pmb { \bar { z } } . { s } }$needs to be able to satisfy the
conditions ot the trace () in star multiplication with$e . { z } \star$
$\cdots \star \mathbf { e } . { \overline { z } }$ the only HS generator that contributes is$V . { x - 1 } ^ { \star }$
generator.When we gauge the field$A . { \bar { \mu }. s }$ $d \bar { z }$component
appears in$. { t i A }$ while $A . { A d S }$and$\lfloor A . { A d S : A } \rfloor.$ do not have
$d \bar { z }$component.The$\bar { \mathbf { A } } . { \mathbf { z } . { \mathbf { n } } }$has$d \bar { z }$component that comes
from $A . { A d S }$ part and it is$e : ^ { p } V . { - 1 } ^ { 2 } e l \overline { z }$ This however will
not appear with the right number of derivatives on$\Lambda$
Since we have chosen$\cdot$to be chiral and$\Lambda = 0$ that was
the only contribution from$\bar { \bar { \boldsymbol { A } } } . { * }.$ Altogether, we can write
$\phi . { z, \dots, z }$component for the$\partial \Lambda ^ { \{ \pmb { \kappa } \} }$derivative as


We have considered the three-point coupling using
metric-like formation to express the higher spin field and
using the linearised Vasiliev's equations of motion. The
obtained result can also be verified using the alternative methods, for example following the procedure by .
The generalisation of the result to higher point functions
would be non-trivial since in order to compute higher order vertices, one would have to consider perturbations
around the background AdS field with higher spin fields
up to that required higher order.


\subsection*{I. CONCLUSION AND OUTLINE}

\end{document}
